\slajdid{09_1-s030}
\begin{frame}[label=09_1-s030]{\alert{Dinamička otpornost MOS tranzistora}}
\gusttekst
Da bi se u kasnijim analizama dinamičkog režima rada tranzistora izbegao prethodni račun, sa dovoljnom tačnošću, MOS FET tranzistor kada radi u aktivnoj oblasti zamenjuje se dinamičkom otpornošću. \alert{Uzima se u obzir i faktor $\lambda$.}
\begin{center}\begin{tikzpicture}[x=.8cm,y=.65cm,font=\sitneoznake,>=Latex]
\draw(-.9,0)--(-.28,0);\node[left]at(-.9,0){$V_{GS}$};\draw(-.28,-.3)--(-.28,.3);\draw(-.2,-.35)--(-.2,.35);\draw(-.2,.25)--(0,.25)--(0,.7)--(1.5,.7);\draw(-.2,-.25)--(0,-.25)--(0,-.9)node[ground]{};
\draw(1.5,.7)--(1.5,-.05);\draw(1.3,-.05)--(1.7,-.05);\draw(1.3,-.17)--(1.7,-.17);\draw(1.5,-.17)--(1.5,-.9)node[ground]{};\node[right]at(1.75,-.1){$C_L$};\node[above]at(.7,.7){$V_O\ (V_{DD}\to V_{DD}/2)$};\end{tikzpicture}
\hspace{10mm}\begin{tikzpicture}[x=.8cm,y=.6cm,font=\sitneoznake,>=Latex]
\draw[->](-.12,0)--(3.7,0)node[right]{$V_{DS}$};\draw[->](0,-.1)--(0,2.65)node[left]{$I_D$};
\draw(0,0)..controls(.4,1.85)and(.45,1.85)..(.8,1.93)--(3.25,2.38);\node[right]at(3.25,2.38){$V_{GS}$};
\draw[dashed](0,2.36)--(3.15,2.36)--(3.15,0)node[below]{$V_{DD}$};\draw[dashed](0,2.08)--(1.6,2.08)--(1.6,0)node[below]{$V_{DD}/2$};
\draw(3.09,2.3)--(3.21,2.42);\draw(3.09,2.42)--(3.21,2.3);\node[below right]at(3.15,2.36){1};\draw(1.54,2.02)--(1.66,2.14);\draw(1.54,2.14)--(1.66,2.02);\node[below left]at(1.6,2.08){2};\draw[->](2.7,2.28)--(2.3,2.20);
\end{tikzpicture}
\end{center}
Ideja je da se prilikom procene kašnjenja, ne izvode i ne računaju kvadratne jednačine, nego da se tranzistor zameni ekvivalentnom otpornošću kada menja radnu tačku između tačaka 1 i 2. Time bi koristili gotove rezultate $t_p=0.69R_nC_L$ sa dovoljno dobrom tačnošću.\par\bigskip
\alert{Uočiti da će ova aproksimacija praktično važiti samo za tranzistore sa kratkim kanalom kod kojih je $V_{DSsat}$ malo i gde važi $V_{DSsat}<\frac{V_{DD}}2$, odnosno gde tranzistor ostaje u zasićenju sve vreme kada prazni parazitnu kapacitivnost, bez obzira koja je pobuda $V_{Tn}\leq V_{GS}\leq V_{DD}$. \textbf{Tranzistori koji se koriste u savremenim integrisanim kolima jesu tranzistori sa kratkim kanalom.}}
\end{frame}
